\section{The automorphic range and the domains of orbital integration}
\label{sec:period-range-v59}
Fix the original datum \(D\), its Hida setting \(X\), a supplied
\(Y:\mathrm{Full}(X)\), and a word \(u\). Let
\(G=G_\infty\times G_f\) carry its specified product Haar measure
\(\mu\), let \(\Gamma\) be its rational subgroup, and let
\(F=X.B.\mathrm{fund}\). Write \(m=\mu|_F\), \(K=X.K_2\) and
\(\chi=X.\mathrm{chi2}\). The Hilbert space is
\(L^2(G,m)\), with inner product linear in the second argument.
The rational subgroup is countable, its translates of \(F\) cover
\(G\), and distinct translates overlap only on null sets. Both
left and right translations preserve \(\mu\). At the finite places
this uses unimodularity of the specified unitary group, as in the
Haar construction of the predecessor chapters. Write \(S=S_u(F)\)
for the rational seed stabilizer and \(\Phi=\Phi_u(T)\) for the
scalar seed of an arbitrary original column \(T\).

% BEGIN QPR1
\paragraph{QPR1. Identifying the weak convolution on the quotient.}
Let \((H,\nu)\) be a sigma-finite parameter space, \(i:H\to G\)
measurable, and \(k\in L^1(H,\nu)\). Let \(R\) be the given right
regular action on \(L^2(G,m)\). Suppose a bounded operator \(A\)
satisfies the original weak formula
\[
 \langle a,Ab\rangle=\int_H k(h)\langle a,R(i(h))b\rangle\,d\nu(h).
\]
For a measurable representative \(b_0\) of \(b\) and a measurable
selector \(r(z)\in\Gamma z\cap F\), a representative of \(Ab\) is
\[
 x\longmapsto \int_H k(h)b_0(r(xi(h)))\,d\nu(h)
 \qquad(m\text{-almost every }x).
\]
\begin{proof}
A selector is obtained by enumerating \(\Gamma\) and choosing the
first translate taking \(z\) into \(F\). For each fixed \(g\),
\(x\mapsto r(xg)\) preserves \(m\) as a measure on the chosen
domain. Indeed, \(Fg\) is another fundamental domain by right
invariance. Partition it according to the translate of \(F\) it
meets, then translate each piece back on the left. Boundaries where
the translate is ambiguous are null by countability. The translated
pieces cover \(F\) up to null sets and preserve measure. This also
shows that the supplied regular representative agrees almost
everywhere with \(b_0(r(xg))\), independently of null changes in
\(b_0\) and of the selector.

Choose a measurable version of \(k\). For any measurable \(E\)
of finite \(m\)-measure, Cauchy--Schwarz and the preceding calculation
give
\[
 \int_H\!\int_E |k(h)b_0(r(xi(h)))|\,dm(x)d\nu(h)
 \le \|k\|_1\,m(E)^{1/2}\|b\|_2.
\]
The integrand is jointly measurable. This is a product-space bound,
so Fubini applies before any exceptional sets depending on \(h\)
are combined. Test the weak formula with \(a=\mathbf1_E\). The
result says exactly that \(Ab\) and the displayed scalar integral
have the same integral on \(E\). Both are integrable there. A
sigma-finite exhaustion and testing the positive and negative parts
of the real and imaginary differences give equality almost everywhere.
Replacing the measurable version of \(k\) by its original version
leaves these integrals unchanged. If \(b_0\) is left
\(\Gamma\)-invariant at every point, then \(b_0(r(z))=b_0(z)\)
pointwise. Thus scalar reproduction almost everywhere for the
unrestricted translates implies \(Ab=b\) on the actual quotient.
\end{proof}
% END QPR1

% BEGIN QPR2
\paragraph{QPR2. Archimedean reproduction survives periodization.}
Assume the original \(\mathrm{FoldOK}(X,u,Y.hK)\) and
\(\varphi=X.\mathrm{ph}\in L^1(G_\infty,\mu_\infty)\).
For every permitted transversal \(L\) of \(S\backslash\Gamma\),
write \(P(x)=\sum_{d\in L}\Phi(dx)\). Then
\[
 \int_{G_\infty}\varphi(z)P(x(z,1))\,d\mu_\infty(z)=P(x)
 \qquad (\mu\text{-almost every }x).
\]
\begin{proof}
Let \(B(x)=\sum_{d\in L}|\Phi(dx)|\). Reindexing modulo \(S\)
gives left \(\Gamma\)-invariance of both \(P\) and \(B\).
The nonnegative unfolding in PGR12 gives
\(\int_F B\,d\mu=\int_{F_S}|\Phi|\,d\mu\) for the domain supplied
by folding. This last integral is finite: the defining weighted
integral for a compact \(\Omega_0\ni1\) is finite, and its
multiplicity is at least one everywhere (use the identity rational
element). For every \(z\), right invariance on the quotient now gives
\(\int_F B(x(z,1))\,d\mu(x)=\int_F B\,d\mu\).
Consequently
\[
 \int_F\int_{G_\infty}|\varphi(z)|B(x(z,1))\,d\mu_\infty(z)\,d\mu(x)
 =\|\varphi\|_1\int_F B\,d\mu<\infty .
\]
The exceptional set for finiteness of the inner integral is
\(\Gamma\)-invariant; its intersection with \(F\) is null. The
countably many translates of this intersection cover the exceptional
set in \(G\). For every remaining \(x\), the sum of the integrals of
the absolute values of the individual translated seed terms is finite.
Integrate termwise. Each seed integral gives \(\Phi(dx)\) by the
archimedean seed identity of PGR7: on its support a fixed finite
phase multiplies a continuous linear evaluation of the vector
reproduction law for \(W_u\), and outside it both sides vanish.
Summing gives the stated equality.
\end{proof}
% END QPR2

% BEGIN QPR3
\paragraph{QPR3. The original range implication \texttt{InHHC}.}
Under exactly the two hypotheses of QPR2 the actual Hilbert period
\(p=\mathrm{PuFun}(X,u,Y.hK,L,T)\) satisfies
\[
       X.B.\mathrm{Re}\bigl(X.\mathrm{Rph}(p)\bigr)=p .
\]
Together with the pointwise absolute convergence, independence of
transversal as a function, and square integrability of the same
series, this supplies every field of \(\mathrm{InHH}\) at every
column and every permitted transversal.
\begin{proof}
PGR11--PGR12 applied to the compact set and multiplicity bound in
\(\mathrm{FoldOK}\) give the three stated scalar fields. In particular
the defining \(L^2\)-branch of \(\mathrm{PuFun}\) is \([P]\).
The countable sum is measurable and pointwise left
\(\Gamma\)-invariant. Apply QPR1 with \(i(z)=(z,1)\), the weak
formula of \(X.\mathrm{Rph}\), and QPR2. It gives
\(X.\mathrm{Rph}(p)=p\). At the finite places the source kernel is
\[
 e(k)=\mathbf1_K(k)\,\overline{\chi(k)}/\mu_f(K),
 \qquad 0<\mu_f(K)<\infty .
\]
Compact openness and positivity are the existing level properties
(ASG14). PGR7 gives \(P(x(1,k))=\chi(k)P(x)\) on \(K\), since each
seed has that covariance. Thus the finite integrand is exactly
\(\mathbf1_K(k)P(x)/\mu_f(K)\), for every \(x\). Its integral is
\(P(x)\). Apply QPR1 with \(i(k)=(1,k)\) and the weak specification
of \(X.B.\mathrm{Re}\). It too fixes \(p\), which proves the
displayed equation in the indicated order. In the range of the
composite, \(p\) itself is the preimage. With PGR19 this gives the
original linear column map into \(\mathrm{HHX}\). At two words on the
same \(X,Y\), a single hypothesis on \(X.\mathrm{ph}\) and folding
at each word now supply both range inputs of PGR18.
\end{proof}
% END QPR3

% BEGIN QPR4
\paragraph{QPR4. Strict orbital domains for every original term.}
For arbitrary \(X,Y,u,u'\) as above, every term \(t\) in the source
has a Borel fundamental domain in \(G_\infty\) meeting each orbit of
\(\mathrm{ROinf}(t)\) exactly once. This proves \texttt{FundExistsC}.
\begin{proof}
Write \(\Lambda,\Lambda'\) for the two row lattices and \(g\) for the
rational linear isometry of \(t\). The group in question consists of
archimedean components of rational \(n\) in the integral flag
unipotents of \(\Lambda\) such that \(g^{-1}ng\) is in those of
\(\Lambda'\). Rational coordinates respect multiplication and inverse.
The flag conditions define a subgroup: the induced action on each
specified successive quotient is the identity, and the row lattice is
preserved with equality. Hence these two requirements define an
intersection of subgroups, including for a degenerate term.

The row property in \(Y\) supplies a finitely generated full
\(\mathcal O_E\)-lattice \(\Lambda\). In the fixed coordinates of
\(X\), pick finitely many generators \(s\), spanning the ambient
\(E\)-space. One nonzero integer \(c\) clears the denominators of
all their coordinates and hence of the whole lattice. If \(M\) is the
matrix of \(n\), all coordinates of \(c(Ms-s)\) are algebraic
integers. In a neighborhood of the identity all their archimedean
absolute values are less than one, simultaneously for every chosen
generator and coordinate. Indeed at place \(v\) the matrix appearing
in the infinite group is \(P_v^{-1}\iota_v(M)P_v\); conjugating back
by this fixed \(P_v\) gives continuous displacement coordinates
vanishing at one. The chosen embedding over each real place of the
totally real base also controls its complex conjugate. A nonzero
algebraic integer has nonzero integral norm, so cannot have every
conjugate of modulus less than one. All the displacements vanish.
The generators span, so \(M=1\). This isolates the identity and
proves discreteness of the actual subgroup \(R_t\).

Choose an open \(V\) about one with \(VV^{-1}\cap R_t=\{1\}\).
Second countability supplies a countable cover \(Vx_j\) of the
group. Each \(Vx_j\) contains at most one point of each left orbit.
Take
\[
 F_t=\bigcup_j\left(Vx_j\setminus\bigcup_{i<j} R_t Vx_i\right).
\]
The inner unions are open, hence \(F_t\) is Borel. An orbit has a
first index at which it meets the cover, and its unique point there
belongs to \(F_t\). This proves coverage and strict uniqueness. It
also shows that the source's chosen \(\mathrm{FundOf}(t)\) selects
a valid domain. No covolume or positivity of an orbital integral is
needed here.
\end{proof}
% END QPR4

% BEGIN QPR5
\paragraph{QPR5. A finite bound on the multiplicity after unipotent unfolding.}
Fix a left row \(\Lambda\). Write \(N_\Lambda\) for its rational
integral flag unipotents, identified with their infinite components,
and \(N'_{\Lambda'}\) for those of a right row. For each term with
this left row let \(r_t\) denote its chosen rational return and
\(R_t=N_\Lambda\cap r_t N'_{\Lambda'}r_t^{-1}\).
Choose representatives \(m\) of \(R_t\backslash N_\Lambda\).
The map
\[
 (t,m)\longmapsto S_{u'}(F)\,r_t^{-1}m
          \quad\text{in } S_{u'}(F)\backslash\Gamma
\]
has uniformly bounded finite fibers for these fixed data. The bound
may depend on the setting and the word.
\begin{proof}
Write \(S'=S_{u'}(F)\), \(E_0\) for the span of the isotropic
frame of the right word, and \(z\) for its negative vector. The
space \(V_0=E_0^\perp/E_0\) is nondegenerate. The rational group
\(S'\) preserves the flag, acts by the identity on its successive
lines, and fixes the class of \(z\) in \(V_0\): in the unipotent--Levi
splitting of PGR2 the Levi fixes the framed vectors, while the
unipotent translations of \(z\) are in \(E_0\). Thus \(S'\) acts
by isometries on \(W=\bar z^\perp\subset V_0\). This space is
positive definite at all real places. Each of the \(r_S\) pairs of
dual isotropic vectors removes a hyperbolic plane, and the orthogonal
negative vector removes the last negative direction of the given
signature. The projection along \(E\bar z\) onto \(W\) is rational
since \(h(z,z)\ne0\).

For a fixed right lattice \(L\), project \(L\cap E_0^\perp\) first
to \(V_0\), then to \(W\). Its image is a full
\(\mathcal O_E\)-lattice. Indeed the intersection is finitely
generated by Noetherianity and spans \(E_0^\perp\) by clearing
denominators; both projections are rational surjections. This is
the source's \(L_W\) in the representative orthogonal also to the
dual frame. Elements of \(S'\) stabilizing \(L\) preserve this
image in both directions. Its isometry group is finite: under the
positive definite rational trace form, images of each member of a
finite spanning lattice basis have prescribed length and lie in a
discrete \(\mathbb Z\)-lattice in real space. There are finitely many
possibilities. An element in the kernel on \(W\) acts identically
on \(V_0=W\oplus E\bar z\). Together with the flag and lattice
conditions this is exactly the defining condition for membership in
\(N'_L\). Conversely these flag unipotents act trivially there.
It follows that for every pair \(L_1,L_2\) of right rows, the maps
in \(S'\) carrying \(L_2\) to \(L_1\) have finitely many classes
modulo right multiplication by \(N'_{L_2}\): if nonempty, this set
is a translate of the stabilizer of \(L_2\). Choose one representative
per class and sum the resulting cardinalities over the finite set of
ordered row pairs to get \(M\).

Within a nonempty fiber choose \((t_0,m_0)\) and put
\(g_i=r_{t_i}^{-1}m_i\) for any of its indices. Then
\(g_0g_i^{-1}\in S'\) carries the right lattice of \(t_i\) to that
of \(t_0\). Label \(i\) by its right lattice and the selected
representative of this map. Equal labels for \(i,j\) imply
\(g_i g_j^{-1}\in N'_{\Lambda'}\). Hence
\(r_{t_i}^{-1}m_i m_j^{-1}r_{t_j}\in N'_{\Lambda'}\).
The two returns represent the same double-unipotent term, so the
terms and then their chosen representatives agree. Now
\(m_i m_j^{-1}\in R_t\); uniqueness in the chosen transversal
gives \(m_i=m_j\). The labels inject this fiber into a set of size
at most \(M\), as required.
\end{proof}
% END QPR5

% BEGIN QPR6
\paragraph{QPR6. Absolute summability over all orbital terms from finite support.}
Retain the two original folding hypotheses. The additional finite-place
input in this claim is the following precise support assertion. Put
\(a_\Lambda=h_{\mathrm{base}(\Lambda)}x_u\). For every original
term \(t\) with left lattice \(\Lambda\), there exist
\(h'\in H_{u',f}\) and \(k\in K\) such that
\[
           r_{t,f}^{-1}a_\Lambda=h'x_{u'}k.                 \tag{QPR.1}
\]
This is the conclusion of \texttt{SlabMajPt.term\_supp}. Then
\[
  \sum_t J_t<\infty,\qquad
  J_t=\int_{\mathrm{FundOf}(t)}
  \|W_u(x)\|\,\|W_{u'}(r_{t,\infty}^{-1}x)\|\,d\mu_\infty(x).
\]
The sum runs over the entire original term type. The archimedean
\(L^1\) hypothesis on \(X.\mathrm{ph}\) in QPR2 is unnecessary here.
\begin{proof}
For each word choose an orthonormal basis of its wave space. For
every finite Levi column class and basis vector choose by PGR6 the
column whose Riesz coordinate is that vector on that class and zero
on the others. This gives a finite list \(T_i\), and likewise
\(T'_j\). At a chosen representative the sum of the absolute
evaluations bounds the vector norm since \(\|v\|_2\le\|v\|_1\) in
orthonormal coordinates. Replacing the representative acts on the
vector by a unitary rational Levi matrix and multiplies by a scalar
of modulus one. Thus this bound holds for any finite Levi component.
The extra level factor also has modulus one. Apply it in the left
column at \(a_\Lambda\) and in the right column supplied by (QPR.1).
Writing \(\Phi_i,\Phi'_j\) for the two lists of seeds bounds the
integrand of \(J_t\) by the sum over \(i,j\) of
\[
 H_t(x)=|\Phi_i(x,a_\Lambda)|\,
        |\Phi'_j(r_t^{-1}(x,a_\Lambda))|.
\]
Take a strict domain \(U_\Lambda\) for \(N_\Lambda\); QPR4's
construction applies directly to this group. For every \(m\in
N_\Lambda\), its unique rational adelic lift \(\tilde m\) belongs
to \(S_u(F)\) and satisfies \(a_\Lambda^{-1}\tilde m_f a_\Lambda\in K\)
by PGR9 at the base column. Moving this finite factor to the right
changes neither absolute value. On the first seed, left invariance
removes \(\tilde m\); on the second seed its new argument is
\(r_t^{-1}\tilde m(x,a_\Lambda)\). In particular \(H_t\) is
\(R_t\)-invariant since for \(m\in R_t\) the conjugated element is
a right-row rational unipotent. Tiling an \(R_t\)-domain by
\(m U_\Lambda\) over the transversal in QPR5 therefore gives, by
Tonelli and left invariance, the bound
\[
 \sum_{t:\,\mathrm{left}(t)=\Lambda}
      \int_{\mathrm{FundOf}(t)} H_t
 \le M\!\int_{U_\Lambda}\!|\Phi_i(x,a_\Lambda)|
        \sum_{d\in S'\backslash\Gamma}
             |\Phi'_j(d(x,a_\Lambda))|\,d\mu_\infty(x).
\]
Here \(d\) denotes any chosen representative. Each absolute
summand is constant on its coset. Changing from the tiled domain
to the chosen strict domain uses the intersection partition under
the countable group \(R_t\); all integrands so far are nonnegative.
The right side is finite by PGR13's absolute slab estimate for this
pair of columns, which itself uses just the two fold bounds in this
application. Sum over the finitely many left rows and basis columns.
This proves the displayed all-term bound, also when either basis is
empty.
\end{proof}
% END QPR6

% BEGIN QPR7
\paragraph{QPR7. What the all-term bound supplies on the actual pair.}
Under the hypotheses of QPR6 the original \texttt{OOIntAll} and all
three clauses of \texttt{AbsConvT} hold on \(X,Y,u,u'\).
\begin{proof}
The rank-one integrand sends \(v\) to
\(\langle W_{u'}(r_t^{-1}x),v\rangle W_u(x)\), with \(r_t\)
here taken at infinity. Its norm is the product of the two vector
norms, including if either vector vanishes. The waves are continuous
in finite dimension, so this is a strongly measurable operator
integrand. Each \(J_t\) is finite, giving Bochner integrability over
the stipulated domain and \(\|\mathrm{OO}_t\|\le J_t\).

Put \(b=|\mathrm{scl}(Y.FN_u,Y.FN_{u'})\,X.\mathrm{volK2}|\).
Every original block is zero or a unit phase times that scalar
normalization times \(\mathrm{OO}_t\). The finite block assembly
followed by the source Hilbert--Schmidt embedding has bounded norm
as a linear map. Thus \(\|\mathrm{op}_t\|\le C b J_t\) with a
constant for this fixed pair. The separate vector has just one
row-pair coordinate, the Hilbert--Schmidt image of the normalized
orbital amplitude multiplied by the weight
\(\sqrt{w(\Lambda)w'(\Lambda')}\). The additive character in that
amplitude has modulus one by PGR1. Thus
\[
 \|\mathrm{vec}_t\|
  \le C' b\left(\sum_{\Lambda,\Lambda'}
                  \sqrt{w(\Lambda)w'(\Lambda')}\right)J_t .
\]
Both constants may depend on the finite spaces; summation proves
the two different norm conditions. Finally \(\mathrm{Gam1}(Y.TR)\)
is the product over the supplied places of
\( |\mathrm{GL}_k(\mathbb F_v)|/|U_k(\mathbb F_v)| \).
At the identity the finite group sum defining the local factor
counts the whole group with the indicated unipotent denominator.
Both group orders are positive, including at rank zero, so their
product is strictly positive. This proves the third clause and the
claimed original outputs. Neither normalizer has been divided out
in this argument.
\end{proof}
% END QPR7

% BEGIN QPR8
\paragraph{QPR8. A finite Levi lift on the given frame supplies every return.}
Equation (QPR.1) holds for every original term. Consequently the
all-term estimate \texttt{OOLintSumC}, the original \texttt{OOIntAllC}
and full \texttt{AbsConvC} hold, with their printed hypotheses.

\begin{proof}
Denote the level modulus by \(g\), the integral finite adelic
coefficient ring by \(A\), the Hermitian involution by a bar, and
the reference lattice by \(L_0\). We first construct the finite Levi
correction for an arbitrary word \(w\). Let \(\epsilon\) be the
integral idempotent of the base finite adeles with component one
above \(g\), zero elsewhere, extended to \(A\). Thus
\(\bar\epsilon=\epsilon\) and \(1-\epsilon\in gA\). On the
\(\epsilon\)-part each \(1+gs\), \(s\in A\), has an integral
inverse: this can be checked in each local integer ring above \(g\).
Equivalently, take the relative quadratic norm
\((1+gs)(1+g\bar s)=1+g(s+\bar s+gs\bar s)\).
The parenthesis is integral in the base finite adeles by integrality
of trace and norm. The base inverse has integral components above
\(g\), and can be extended by zero off those places. Multiplication
by \(1+g\bar s\) gives the required inverse on that summand. The
same conclusion for a matrix congruent to the identity on this
summand follows by its determinant and adjugate.

Use the source's fixed vectors \(U=(e_a,f_a,x_w^{-1}z_w)\) in its
finite reference coordinates (the vectors \(e_a,f_a\) are fixed by
\(x_w\)). On the \(\epsilon\)-part the vectors are integral and
congruent modulo \(gL_0\) to the reference vectors
\(U_0=(e_a^0,f_a^0,z^0)\). These are precisely the transport
congruences above \(g\) used in PGR3. Their Gram matrix there is
congruent to the direct sum \(Q_0\) of hyperbolic pair matrices
and \((-1)\); in particular \(Q_0^2=1\). Thus the Gram matrix
of \(\epsilon U\) is invertible over \(\epsilon A\) by the preceding
argument applied after multiplication by \(Q_0\). For an integral
vector \(v\), subtract from its \(\epsilon\)-component the vector
in the span of \(\epsilon U\) with coefficients
\(\operatorname{Gram}(\epsilon U)^{-1}(h(\epsilon U_i,v))_i\).
This is the integral orthogonal projection away from those fixed
vectors on that summand. Apply it to the supplied \(bw_j,bwd_i\).
All correction coefficients are zero modulo \(g\): the reference
pairings with \(U_0\) vanish modulo \(g\), and
\(\epsilon U\equiv U_0\) there.

Let \(e_w=X.\mathrm{ew}\) denote the source's split coefficient
idempotent (independent of the word), and put \(\eta=\epsilon e_w\).
It satisfies \(\eta\bar\eta=0\) and \(\eta\equiv e_w\bmod gA\).
Take the \(\eta\)-parts of the projected \(bw\)'s and the
\(\bar\eta\)-parts of the projected \(bwd\)'s. Their cross-pairing
matrix \(H=(h(d'_i,b'_j))\) is the identity modulo \(g\) on the
\(\eta\)-summand by the original dual-frame congruences. Correct
the \(b'\)-columns by \(H^{-1}\) on that summand. The resulting
integral families \(b,d\) are orthogonal to all the actual fixed
vectors, satisfy \(h(d_i,b_j)=\delta_{ij}\eta\), and agree with
the original two families modulo \(gL_0\).
Both same-side pairings are zero by orthogonality of the two
coefficient idempotents. Thus their span splits integrally off
\(L_0\): its two coordinate projections are given by the pairings
against \(d_i,b_j\). Everything outside this span may be fixed.

Let \((s_v)_v\) be any tuple in the source product
\(\prod_v\mathrm{GL}_k(Y.TR.Fq_v)\). Joint surjectivity of the
given reduction maps on \(A\) supplies integral matrices \(C,C_0'\)
reducing to \(s_v,s_v^{-1}\), respectively. Use the actual
\(Y.TR\) throughout. Its kernel condition says that an entry with
all reductions zero has its \(e_w\)-part in \(gA\). Therefore
\(CC_0'\) and \(C_0'C\) are identity modulo \(g\) on the
\(\eta\)-part. The determinant argument there supplies a genuine
integral inverse \(C'\) of \(C\) on this part. Define the
transformation \(l\) to act on the \(b\)-columns by \(C\), on
the \(d\)-columns by \(\overline{C'}^{\,t}\), and by identity on
the orthogonal complement. The exact cross-pairing shows it is
unitary. Both it and its inverse preserve \(L_0\), and it fixes the
specified \(U\). Its reduced matrix on the \(bw\)-quotient has
entries \(\mathrm{res}_v h(bwd_i,l\,bw_j)=(s_v)_{ij}\) by the
congruences of both frames. In particular \(x_w l x_w^{-1}\) fixes
the rational framed vectors exactly and belongs to \(H_{w,f}\).
This constructs the correction without choosing new residue maps or
requiring the two idempotents to exhaust the coefficient ring.

Now write \(a=h_{\mathrm{base}(\Lambda)}x_u\),
\(a'=h_{\mathrm{base}(\Lambda')}x_{u'}\), and \(k_0=a^{-1}r_{t,f}a'\).
The equality of rational row lattices in the carrier gives
\(k_0L_0=L_0\) on completion. On the reference \((e^0,z^0)\)-flag
modulo \(gL_0\), \(k_0\) is upper unitriangular. Indeed the source
admissibility witness supplies one integral upper unitriangular
matrix for the rational frames. Read its column congruences at
places over \(g\) using the lattice identification of the base
columns and the same transport congruences as above (PGR14).
Away from \(g\), an integral displacement is automatically in
\(gL_0\). This proves the first two clauses defining \(K_2\) for
\(k_0\) and also for its inverse by triangular inversion.

For clarity, the residual Levi matrix here is taken on the
\(bw\)-block modulo the span of \(e^0\). Indeed for such an
integral unitary element the images of \(bw_j\) pair to zero modulo
\(g\) with \(e_a^0,z^0\) by applying its inverse to these vectors;
their pairings with \(bw_i\) vanish on the common idempotent side.
In the supplied reference dual coordinates these conditions put
the images in the span of \((e^0,bw)\) modulo \(gL_0\). The
coordinates on the latter block are read by \(h(bwd_i,-)\).
This holds for the inverse too, hence the resulting reduced
matrices \(s_v(k_0)\) are invertible, and composition multiplies
them in order. The reference coordinate identity used here is
\(v\equiv\mathrm{repr}_0(v)\bmod gL_0\) for integral \(v\);
it expands \(v\) in the reference vectors with the corresponding
dual pairings (including the minus sign on the \(z^0\)-coordinate).
Thus no assumption that the rational row vectors form an integral
basis was used.

Apply the preceding lifting construction with \(w=u'\) and
\(s_v=s_v(k_0)^{-1}\). The correction fixes \(U\), so it fixes
the reference \((e^0,z^0)\)-flag modulo \(gL_0\). The product
\(k_0l\) retains those congruences and has identity reduced
\(bw\)-block. The coefficients of \((k_0l-1)bw_j\) on the
\(bw\)-block are supported on \(e_w\); the typed kernel criterion
puts them in \(gA\). The only remaining components modulo \(gL_0\)
are then in the span of \(e^0\). This is stronger than the third
\(K_2\) clause, which also permits earlier \(bw\)'s.
Therefore \(k_0l\in K_2\). With \(h'=h_{\mathrm{base}(\Lambda')}
(x_{u'}l x_{u'}^{-1})\in H_{u',f}\) the identity
\[
        r_{t,f}^{-1}a=h'x_{u'}(k_0l)^{-1}
\]
gives (QPR.1), uniformly in the sense of existence for each
original term. Applying QPR6 and QPR7 at the two supplied fold
bounds proves the assertions.
\end{proof}
% END QPR8
% BEGIN QPR9
\paragraph{QPR9. Stabilizer phases and independence of the strict domain.}
Fix $D,X,Y,u,u'$ and lattices $\Lambda\in Y.\mathrm{what}_2(u)$,
$\Lambda'\in Y.\mathrm{what}_2(u')$. A carrier $r$ consists of a
rational map with its linked adelic point. Its real stabilizer in the
orbital integral is the infinite image of
\[
 R_r=N_{u,\Lambda}\cap rN_{u',\Lambda'}r^{-1}.
\]
The two groups on the right are the actual rational flag unipotents
preserving the respective lattices. Rational links commute with
products and inverses, by matrix multiplication in all embeddings.
They are unique since the coordinates of a rational linear map
specify its adele at every place.

Let $g_*$ be the level modulus and work first modulo $g_*L_0$.
Write $E_{<a}$ for the span of the reference vectors $e_b^0$, $b<a$.
Consider the integral finite isometries $k$ satisfying
$(k-1)e_a^0\in E_{<a}$ and $(k-1)z^0\in\sum_a A e_a^0$
modulo this modulus, with $A$ the integral coefficient ring.
They form a group: on the flag spanned by $(e^0,z^0)$ they are
upper unitriangular modulo the modulus. Multiplication preserves the
congruences, and on this flag the inverse is a finite geometric
polynomial in the strictly triangular displacement.

On this whole group the source character $\chi_N$ is multiplicative.
Its numerator is $\sum_a h(f_a^0,k\,v_a)$, where
$v_a=e_{a+1}^0$ except at the last step where $v_a=z^0$.
The numerator for $k=1$ vanishes modulo $g_*$, by orthogonality
of distinct framed vectors and $h(f_a^0,z^0)\equiv0$. For two
elements $k,l$ use
$kl-1=(k-1)+(l-1)+(k-1)(l-1)$. The vector $(l-1)v_a$
lies in $E_{<a+1}$ modulo $g_*L_0$, including the endpoint $z^0$.
Applying $k-1$ puts it in $E_{<a}$ there; the error stays in
$g_*L_0$ since $k$ is integral. Its pairing with $f_a^0$ therefore
vanishes modulo $g_*$. After dividing by $g_*$ the extra term has
trivial finite additive character by PGR1. We have shown
$\chi_N(1)=1$ and $\chi_N(kl)=\chi_N(k)\chi_N(l)$.
In particular it is invariant under conjugation inside this group.

For $x\in R_r$, put $x'=r^{-1}xr$ and use the two base columns with
transports $a,b$. Then $k=a^{-1}r_f b$ lies in the integral
congruence group just considered, by the base-column congruences
proved in QPR8. That argument uses only the defining lattice
identification and admissible frame of a carrier and so applies to
every $r$ here. Put $d=a^{-1}x_f a$ and $d'=b^{-1}x'_f b$.
PGR9 puts these in $K_2$ and identifies their $\chi_2$ phases with
$\psi_{u,f}(x_f)$ and $\psi_{u',f}(x'_f)$. Their $\chi_N$ phases
have exactly the same values: in the proof of PGR3 the entire
additional $bw$ numerator of $\chi_2$ vanishes modulo $g_*$ on
these conjugated unipotents, also after the column's Levi transport
as in PGR9. Now $d=kd'k^{-1}$. Thus the finite phases coincide,
even when $k$ fails the extra $bw$ conditions for $K_2$.
The rational product formula of PGR9 gives
\[
        \psi_{u,\infty}(x_\infty)
        =\psi_{u',\infty}(r_\infty^{-1}x_\infty r_\infty).
\]
Both phases have absolute value one.

The orbital integrand is the rank-one operator with vectors
$W_u(t)$ and $W_{u'}(r_\infty^{-1}t)$. On translation by
$x_\infty$ the two vectors acquire the two equal phases above;
the representation factor on a flag unipotent is the identity.
The conjugate scalar in the second argument cancels the first scalar,
so the operator is invariant. The rational stabilizer is countable
since it consists of matrices over a number field. Given two strict
measurable domains $\Omega,\Omega'$, partition $\Omega$ by its
intersections with the stabilizer translates of $\Omega'$. Their
inverse translates partition $\Omega'$. Left invariance of the same
Haar measure and integrand invariance identify the integrals on
corresponding pieces. The integral of the norm is preserved by
this partition as well. Hence integrability transfers in either
direction and the Bochner integrals agree whenever integrable
(and agree under the source's zero convention otherwise). With
the domain existence and all-term integrability premises of
\texttt{OOFundIndepC}, specialize $r$ to each term representative
and $\Omega'$ to its selected \texttt{FundOf}. This proves precisely
$\mathrm{OO}(\Omega)=\mathrm{OOt}$ for every permitted $\Omega$.
% END QPR9

% BEGIN QPR10
\paragraph{QPR10. Both amplitudes under a change of carrier.}
With the same $D,X,Y,u,u'$, take arbitrary columns $\rho,\rho'$
and their transports $a=h_\rho x_u$, $b=h_{\rho'}x_{u'}$. Let
$r,s$ be related carriers on those column lattices. The relation
supplies $s=n r n'$ on rational maps for
$n\in N_{u,\Lambda_\rho}$ and $n'\in N_{u',\Lambda_{\rho'}}$.
A rational isometry determines its adelic image by acting on the
supplied rational basis at all places; it preserves the Hermitian
form there, and its finite matrix is integral outside finitely many
places, so it gives the required rational link. Its uniqueness makes
$s=n r n'$ true on the actual adelic points as well. Write
\[
 d=a^{-1}n_f a,\qquad k_r=a^{-1}r_f b,\qquad d'=b^{-1}n'_f b .
\]
Then $k_s=d k_r d'$. Both $d,d'$ lie in $K_2$ by PGR9.
Consequently $k_s\in K_2$ if and only if $k_r\in K_2$, the reverse
implication following by multiplication by their inverses. This
establishes the support clause for every related pair.

Choose any two strict measurable domains $\Omega_r,\Omega_s$ for
the two real stabilizers. In the notation of QPR9,
$R_s=nR_r n^{-1}$ because $n,n'$ lie in the indicated lattice
unipotent groups. Hence $n_\infty^{-1}\Omega_s$ is another strict
domain for $R_r$. Substitute $t=n_\infty z$ in the integral over
$\Omega_s$. The first vector acquires $\psi_{u,\infty}(n_\infty)$.
In the second vector the new argument is
$n_\infty'^{-1}r_\infty^{-1}z$; its coefficient is the inverse
of $\psi_{u',\infty}(n'_\infty)$. Since the latter is unitary,
conjugate linearity in that vector converts the inverse to
$\psi_{u',\infty}(n'_\infty)$. Applying the domain comparison
from the proof of QPR9 on $R_r$ gives
\[
 \mathrm{OO}(s,\Omega_s)=
 \psi_{u,\infty}(n_\infty)\psi_{u',\infty}(n'_\infty)
                         \,\mathrm{OO}(r,\Omega_r).
\]
This applies to the source's prescribed integrability premise on
$\Omega_r$ and transfers integrability to the translated domain.

If $k_r\in K_2$, for either $\xi=\chi_2$ or $\xi=\chi_N$ we have
\[
 \xi(k_s)=\psi_{u,f}(n_f)\,\xi(k_r)\,\psi_{u',f}(n'_f).
\]
For $\chi_2$ use the original level character law on $K_2$; for
$\chi_N$ use QPR9 since $K_2$ satisfies the required integral
flag congruences. PGR3 and PGR9 give both outer factors for the
actual columns. Multiplying the two displays cancels each rational
infinite phase against its own finite phase by PGR9. It follows
that $\xi(k_r)\mathrm{OO}(r,\Omega_r)=
\xi(k_s)\mathrm{OO}(s,\Omega_s)$ for each of the two characters.
This proves the universal \texttt{AmpWellDefC} implication with
the original \texttt{OOIntAll}, every pair of columns, and both
domain choices.
% END QPR10

\section{Folding on the original arithmetic quotient}
\label{sec:folding-v059}
Use the notation of QPR1 and of PGR7--PGR13. Here $F$ is the totally
real field, $E/F$ the CM field, $n$ the matrix size, and $q_v$ the
negative index at $v$. Write $S_\infty=N_\infty H_\infty$ for the
real stabilizer, $K=X.K_2$ for the finite level, and $W=W_u$.
Set
\[
  L(g)=\prod_{v\mid\infty}\|g_v\|_{\rm op},\quad
  \nu(x)=\inf_{s\in S_\infty}L(sx),\quad
  B=4n^2\sqrt{2\max_v q_v}.
\]
These are the norm, quotient weight and exponent of the original
$\mathrm{WuDecay}$. Its two premises for the fixed word supply a
\emph{compactly supported} reproducing kernel $\alpha$ and finiteness
of $\int \|W(x)\|\nu(x)^B\,d(X.\nu_Q)(x)$.

% BEGIN FLD1
\paragraph{FLD1. Counting rational translates with the prescribed exponent.}
\label{claim:FLD1}
For all compact $\Omega,C_0\subset G_\infty\times G_f$ there is a
finite constant $A$ such that
\[
 m_\Omega(c(g,1))\le A L(g)^B\qquad(c\in C_0,\ g\in G_\infty).
\]
Here $m_\Omega$ is the supremum over the second argument in (PGR.22).

\begin{proof}
Put $x=c(g,1)$. Fix that second argument $y$. If the counted set
is nonempty choose $\gamma_0$ in it. Right division by $\gamma_0$
injects it into $\Gamma\cap x\Omega\Omega^{-1}x^{-1}$, independently
of $y$. The finite components of the latter elements lie in one fixed
compact set $C_f$, the projection of $C_0$ and $\Omega$ being compact.
The stabilizer $K_0$ of the original finite reference lattice is open
(it contains the open first principal congruence subgroup), so finitely
many translates of $K_0$ cover $C_f$. These matrices have one common
denominator. Explicitly, if $\widehat L$ is the completed reference
lattice and $A_f^0$ the integral coefficient ring, finite generation
and full span give
$d(A_f^0)^n\subset\widehat L\subset d^{-1}(A_f^0)^n$ for one positive
integer $d$: clear the coordinates of the given generators and
of expressions for the coordinate vectors in those generators.
A matrix preserving $\widehat L$ therefore has entries in
$d^{-2}A_f^0$. Clear the finitely many coset representatives as
well: an individual coefficient is a finite sum of pure tensors,
whose number-field and finite adelic factors each have bounded
denominators. For a diagonal rational coefficient $b\in E$, this
implies $D b\in\mathcal O_E$ for a common integer $D$, since it
is integral in every finite completion of $E$ after multiplication
by $D$. The integral adelic tensor ring maps into those local
integer rings. This argument only uses the given congruence topology.

At a real place $v$, use the supplied complex matrix coordinates.
Changing coordinates by the fixed $P_v$, the compact projections
of $C_0,\Omega\Omega^{-1}$ and inversion on $C_0$ only contribute
a constant. In signature coordinates $g_v^{-1}=J_v g_v^*J_v$
with $J_v$ a diagonal isometry, hence
$\|g_v^{-1}\|_{\rm op}=\|g_v\|_{\rm op}\ge1$ ($n>0$).
Each entry of the rational matrix thus has
$|\iota_v(b)|\le C_v\|g_v\|_{\rm op}^2$ with $C_v\ge1$.
Each $v$ selects one member of a complex conjugate pair. The
Minkowski embedding of $D^{-1}\mathcal O_E$ is a full lattice in
$\prod_{v\mid\infty}\mathbb C$. The number of its points in disks
of radii $R_v\ge1$ is bounded by $C\prod_v R_v^2$: translate one
bounded fundamental parallelotope about every such point and compare
its volume with that of the product of disks enlarged by its fixed
coordinate diameters. Taking $n^2$ entries bounds the number of
matrices by $C' L(g)^{4n^2}$. A rational matrix determines both
adelic components, so the same estimate counts group elements.
If some $q_v>0$ then $B\ge4n^2$; if all are zero every $g_v$ is
unitary for the positive form and $L(g)=1$. Both cases give the
claimed exponent, including $B=0$.
\end{proof}
% END FLD1

% BEGIN FLD2
\paragraph{FLD2. The weighted integral on every rational stabilizer domain.}
\label{claim:FLD2}
For every compact $\Omega$, every column $T$, and every original
fundamental domain $F_S$ for the rational stabilizer, one has
\[
             \int_{F_S}m_\Omega(y)|\Phi_u(T)(y)|\,d\mu(y)<\infty
\]
under the original weighted wave hypothesis.

\begin{proof}
Index the finite column classes by $\rho$, write $a_\rho=h_\rho x_u$,
and let $N_\rho$ denote PGR9's integral rational unipotents for the
row lattice $\Lambda_\rho$. We need a compact covering set in the
\emph{real} unipotent group, stronger than existence of a measurable
domain in $G_\infty$. Here is a direct construction. Regard the
Lie algebra of the rational unipotent as a nilpotent $\mathbb Q$-Lie
algebra by restriction of scalars. Its real points are precisely the
product of its points over the infinite places. Exponential and
logarithm are inverse triangular matrix polynomials by PGR8.
Choose a rational basis $X_1,\ldots,X_m$ such that
$[\mathfrak n,X_i]\subset\langle X_1,\ldots,X_{i-1}\rangle$.
Such a basis is obtained successively by choosing a central rational
line in each nonzero nilpotent quotient and taking inverse images.
Scale the basis vectors by positive integers so that both
$\exp(X_i)$ and $\exp(-X_i)$ preserve the given row lattice:
in any $\mathbb Z$-basis of that lattice their finite exponential
polynomials have integral coefficients after clearing the rational
matrix entries and factorials in the nonconstant terms. Thus these
elements belong to $N_\rho$. Reducing a real unipotent by their
integer powers on the left gives a compact set of representatives.
Indeed the ordered exponential coordinates
$\exp(t_1X_1)\cdots\exp(t_mX_m)$ exist uniquely, successively
passing to the last one-dimensional quotient and its kernel.
Left multiplication by $\exp(kX_i)$ adds $k$ to coordinate $i$
and changes only lower coordinates, by the displayed bracket rule.
Reduce coordinates in decreasing order into $[0,1]$. Consequently
$N_\infty=N_\rho C_\rho^N$ for a compact $C_\rho^N\subset N_\infty$.
This includes the zero-dimensional group.

The residual Levi $H_\infty$ fixes both hyperbolic frames and the
negative vector. It acts on their orthogonal complement, a positive
definite Hermitian space at every infinite place (the signature
calculation in QPR5), and is a closed subgroup of a product of
compact unitary groups. PGR2 gives
$S_\infty=N_\rho C_\rho$ with $C_\rho=C_\rho^N H_\infty$ compact.
For any $t\in S_\infty$ write $t^{-1}=n c$ with $n\in N_\rho$,
$c\in C_\rho$. Its unique rational lift $\tilde n$ satisfies
$\tilde n_f^{-1}a_\rho\in a_\rho K$ by PGR9. Therefore, for
every $k\in K$,
\[
    (x,a_\rho k)=\tilde n\,
          (c,\tilde n_f^{-1}a_\rho k)\,(t x,1).
\]
Left rational translation preserves $m_\Omega$ by reindexing.
Apply FLD1 with the compact $C_\rho\times a_\rho K$ and take the
infimum over $t$. This gives $m_\Omega(x,a_\rho k)\le A_\rho\nu(x)^B$.
The inequality holds throughout the finite coset, with no claim that
multiplicity is right $K$-invariant. On that coset PGR6--PGR7 also
give $|\Phi_u(T)(x,a_\rho k)|\le\|J_u T(\rho)\|\|W(x)\|$.

Choose a Borel strict domain $U_\rho$ for $N_\rho$ in $G_\infty$ as
in PGR8. We compare
$\int_{U_\rho}\nu(x)^B\|W(x)\|\,d\mu_\infty(x)$ with the original
quotient integral. The weight $w(x)=\nu(x)^B\|W(x)\|$ is
$S_\infty$-invariant by reindexing the infimum and the unitary wave
covariance; it is Borel since the infimum of continuous functions
is upper semicontinuous, and $\nu\ge1$. Apply the supplied Weil
formula to $\mathbf1_{U_\rho}w$. For fixed $h\in H_\infty$ and
$x\in G_\infty$, the set $\{n\in N_\infty:n h x\in U_\rho\}$
is a Borel fundamental domain for $N_\rho$ in $N_\infty$.
All such domains have the same measure by the countable intersection
partition. It is finite: the compact covering set $C_\rho^N$ above
meets every orbit, and a measurable strict domain can be moved into
it by partitioning according to the first hit in an enumeration
of $N_\rho$. These translated pieces are disjoint.

The measures supplied as $X.\mu_N$ and $X.\mu_H$ restrict to genuine
finite-on-compact Haar measures. To see this even from their printed
weaker description, each is nonzero, sigma-finite, supported on its
closed subgroup, and invariant under its left translations. Both
subgroups are unimodular (compact for $H_\infty$, nilpotent for
$N_\infty$, whose adjoint determinants are one). If $m$ is this
restricted measure and $m_0$ a usual Haar measure, Fubini applied to
$f(y)h(zy)$ under $dm(y)dm_0(z)$ gives
$m(f)m_0(h)=m(h)m_0(f)$, by translation invariance of $m$ and
unimodularity of $m_0$. Choose strictly positive measurable $f$
with finite integral under both, using their sigma-finite exhaustions.
Thus $m$ is a finite positive multiple of $m_0$. In particular the
measure of $H_\infty$ is finite. The original Weil identity now gives
\[
 \int_{U_\rho}w\,d\mu_\infty
   =\operatorname{covol}_{\mu_N}(N_\rho)\,
             \mu_H(H_\infty)\int w\,d(X.\nu_Q)<\infty .
\]
This uses exactly its given normalization and the weighted wave premise.
For arbitrary compact $\Omega$, $m_\Omega$ is Borel by the countable
tuple and compact projection argument of PGR11; that measurability
argument applies also when $1\notin\Omega$. By the PGR10 support
partition, integration of the nonnegative
invariant integrand over $F_S$ equals the sum over the cells
$U_\rho\times a_\rho K$ by left translations. For this step retain
the integral over each whole finite coset, since $m_\Omega$ need
not be right $K$-invariant. Its uniform bound on that coset costs
just the finite factor $\mu_f(K)$. Sum over the finitely many $\rho$.
This proves the claim on every allowed $F_S$.
\end{proof}
% END FLD2

% BEGIN FLD3
\paragraph{FLD3. The original folding implication.}
\label{claim:FLD3}
For every original $D,X,Y,u$, level normalization and
$\mathrm{WuDecay}(X)$ imply $\mathrm{FoldOK}(X,u,Y.hK)$.

\begin{proof}
PGR8 gives a domain for $S_u(F)$ in $G$: the rational subgroup is
discrete and countable and the first-hit construction gives a Borel
strict domain, hence an admissible measure fundamental domain.
For the two uniform quantifiers of the averaging inequality, take
the compact $\Omega_i\ni1$, kernel $\alpha$ and bound $a_i$ supplied
by $\mathrm{WuReprod}$ for this word. Let
$\Omega_0=\Omega_i\times K$ and $a=a_i/\mu_f(K)$; ASG14 gives
$0<\mu_f(K)<\infty$. With the finite argument of the seed held fixed,
it is zero outside its support and on support is a fixed complex
linear evaluation of the archimedean wave, times a unit scalar.
Push that functional through the supplied Bochner reproduction
identity. Taking absolute values bounds the seed by $a_i$ times
the integral of its absolute value over right translations by
$\Omega_i$. Right multiplication of its finite argument by $K$
preserves absolute values, by PGR7, so averaging there rewrites
the bound as
\[
      |\Phi(y)|\le a\int_{\Omega_0}|\Phi(y\omega)|\,d\mu(\omega).
\]
The same $a$ and $\Omega_0$ work for every $T$, and the inequality
persists for every compact $\Omega\supset\Omega_0$ by monotonicity.
Measurability and rational left invariance are PGR7; FLD2 gives the
multiplicity integral for every compact $\Omega$ and every choice
of domain. These are exactly the two clauses of $\mathrm{FoldOK}$.
In particular QPR3 applies at each word under the additional
$L^1$ condition on $X.\mathrm{ph}$, and QPR7--QPR8 apply to a
pair on this same $X,Y$. The compact averaging kernel used here
comes from $\mathrm{WuDecay}$ and need not equal $X.\mathrm{ph}$.
\end{proof}
% END FLD3


\subsection{The Gram expansion for a tied pair of words}
% BEGIN QPR11
\paragraph{QPR11. The tied period kernel and the absolute convergence of its terms.}
Fix any original datum $D$, a setting $X$, an object $Y:\mathrm{Full}(X)$,
and an ordered pair of words $u,u'$, allowing $u=u'$. Assume the
original $\mathrm{ConcTie}(X,Y)$. Write $A_t$ for the normalized
orbital operator of PGR18, $P_w$ for the Poincar\'e series attached
to a column of word $w$, and $J_w$ for the original antilinear
identification of columns. For each word $w$ use the original real-valued masses and set
\[
 c_w=\frac{X.\mathrm{volK2}\,(\mu_{N,w}(Y.FN(w)))_{\mathbb R}
                         (\mu_{H,w}(H_{w,\infty}))_{\mathbb R}}
                 {\dim_{\mathbb C}V_w},\qquad
 s_{u,u'}=\bigl(\sqrt{c_uc_{u'}}\bigr)^{-1}.
\]
These are the source's total division and inverse, with $0^{-1}=0$. Then there is
a continuous linear operator $B_{u,u'}$ such that
\[
  \sum_t A_t = B_{u,u'},\qquad
  s_{u,u'}\langle P_uT,P_{u'}T'\rangle
       =\langle B_{u,u'}J_{u'}T',J_uT\rangle
       \quad\text{for all }T,T'.
\]
The first equality holds for the net of all finite partial sums.
The same terms have summable operator norms, the auxiliary vectors
have summable norms, and the identity factor is positive. These
are the original \texttt{KExpTied} and \texttt{AbsConvTied} conclusions.

\begin{proof}
Take the level normalization and weighted wave decay from this
$\mathrm{ConcTie}$. FLD3 applied twice gives folding on the two
specified words on $X$ with the same finite character $Y.hK$.
Its compact reproducing kernels may depend on the word. The
integrable archimedean kernel $X.\mathrm{ph}$ used by the two
Hilbert operators is common: its integrability is another field
of this very tie. QPR3 therefore puts both Poincar\'e series in
the required ranges for the rational representatives $Y.Dt(u)$
and $Y.Dt(u')$.

Use QPR4 for strict domains on this pair. QPR5--QPR8 now give
integrability on the selected domains, both norm sums and positivity
on the original $Y.TR$. This is precisely $\mathrm{AbsConv}$ for
$Y.tO(u,u')$. QPR9 identifies the orbital integrals on all other
allowed domains, and QPR10 identifies the amplitudes for all related
carriers of the two columns. These are the remaining hypotheses of
PGR18. Use it on $X.Fu(u),X.Fu(u')$, the two column systems
$Y.cs(u),Y.cs(u')$, and the given unipotent domains
$Y.FN(u),Y.FN(u')$. Its absolutely convergent operator sum and
sesquilinear pairing give the display with the original constants.
None of these selections depends on changing $X$ or replacing $Y$.
In particular, the two inputs called \texttt{KExpTied} and
\texttt{AbsConvTied} in the diagonal and distinct-word bounds are
available simultaneously. The choice of such data across a family,
its comparison isometries, and the other estimates in those bounds
remain separate mathematical tasks.
\end{proof}
% END QPR11
